% exp-fips203-mlkem-pke-stage1-encode-vec 纯向量 Stage1 实现方案
% 编译：bash ../../../../scripts/xelatex-clean.sh exp-fips203-mlkem-pke-stage1-encode-vec-实现方案-customspec.tex
\documentclass[UTF8,a4paper,11pt]{ctexart}
\usepackage{amsmath,amssymb}
\usepackage{booktabs}
\usepackage{geometry}
\usepackage{hyperref}
\usepackage{longtable}
\usepackage{array}
\usepackage{seqsplit}
\geometry{margin=2.2cm}
% 表格列内自动换行；长 API 名在 \texttt 中可断行
\newcolumntype{L}[1]{>{\raggedright\arraybackslash}p{#1}}
\newcommand{\apiname}[1]{\texttt{\seqsplit{#1}}}
\hypersetup{colorlinks=true,linkcolor=blue,urlcolor=blue}

\title{exp-fips203-mlkem-pke-stage1-encode-vec\\F203 Stage1 纯向量实现方案（API 全覆盖评估）}
\author{ascendc 工程（预研）}
\date{2026-05-19}

\begin{document}
\maketitle

\begin{abstract}
本文档描述 \texttt{examples/incubating/exp-fips203-mlkem-pke-stage1-encode-vec/} 用例：
将 F203 交付件 \texttt{sepolyvec8\_ntt\_f203} 的 Stage1（\texttt{se [8,256] int32} $\rightarrow$ \texttt{mat\_a [16,256] int8}，布局 $[HI(8),LO(8)]$）
在 \textbf{Atlas A2 / Ascend910B} 上以 \textbf{纯 AIV 向量}实现。
API 依据 \textbf{CANN 社区版 9.0.0 Ascend C 算子开发接口参考 01}（PDF 第 2 章 SIMD API）及昇腾社区在线页；
dav\_c220（910B）头文件 \texttt{kernel\_operator\_vec\_*\_impl.h} 用于交叉核对 A2 上 \texttt{Cast} 等组合是否真实可用。
CPU 对拍：\texttt{bash run.sh -r cpu -v Ascend910B4}；\textbf{1 / 2 / 8 个 Vector AI Core 三档}（见 §\ref{sec:aicore}）。
\end{abstract}

\tableofcontents
\newpage

\section{数学语义}

\subsection{输入输出}
\begin{itemize}
  \item 输入 \texttt{se\_polyvec\_gm}：形状 $[8,256]$，\texttt{int32}。Host/Python 保证每条系数 $v \in [0,q)$，$q=3329$。
  \item 输出 \texttt{mat\_a\_gm}：形状 $[16,256]$，\texttt{int8}。行 $0..7$ 为 $HI$，行 $8..15$ 为 $LO$。
  \item \textbf{本阶段不做 $\bmod\,q$}（与交付 ONNX 一致；Python \texttt{validate\_se\_polyvec\_in\_zq} 保证前提）。
\end{itemize}

\subsection{逐系数公式（对齐 MlkemEncodeToLimb6 / 交付 Floor-Sub 在 $v\ge0$ 时的等价）}
对每条 poly、每个系数 $v$：
\begin{align}
  hi &= \lfloor v / 64 \rfloor = v \gg 6 \\
  lo &= v - hi \cdot 64 \quad (\text{在 } v\in[0,q) \text{ 时等价于 } v \bmod 64)
\end{align}
\textbf{不采用} $lo = v \land 63$ 的独立 \texttt{And} 路径作为最终实现：CPU 仿真上曾出现 \texttt{lo} 错误并导致 \texttt{Cast} 饱和到 $127$（见 §\ref{sec:rejected}）。

\subsection{Golden}
\texttt{mlkem\_ref.encode\_mat\_a(se)}：\texttt{hi\_i8 = (se>>6).astype(int8)}，\texttt{lo\_i8 = (se\&63).astype(int8)}，沿 axis=0 拼接。

\section{向量实现规划}

\subsection{AI Core 规模与验证档位}
\label{sec:aicore}

\begin{longtable}{@{}L{2.8cm}L{10.5cm}@{}}
\toprule
项 & 本用例约定 \\
\midrule
核类型 & \texttt{KERNEL\_TYPE\_AIV\_ONLY}：仅使用 \textbf{Vector AI Core}（AIV），\textbf{不使用} Cube AI Core（AIC）。 \\
计划 AI Core 数 & \textbf{三档}：\texttt{aiv=1} $\Rightarrow$ \texttt{blockDim=1}；\texttt{aiv=2} $\Rightarrow$ \texttt{blockDim=2}；\texttt{aiv=8} $\Rightarrow$ \texttt{blockDim=8}（\texttt{launch\_profile.h}）。 \\
数据划分（aiv=8） & \texttt{GetBlockIdx()}$=i$：每核 1 poly；读 \texttt{se} 第 $i$ 行；写 \texttt{mat\_a} $HI/LO$ 第 $i$ / $8+i$ 行。 \\
数据划分（aiv=2） & \texttt{blockIdx=0,1}：每核 4 poly；核内 \textbf{for} $p=4\cdot blockIdx .. 4\cdot blockIdx+3$ 串行；\textbf{每 poly 新建} op。 \\
数据划分（aiv=1） & \texttt{GetBlockNum()==1}：核内 \textbf{for} $p=0..7$ 串行；\textbf{每 poly 新建} op（避免 \texttt{TPipe}/\texttt{TQue} 跨 poly 复用）。 \\
多档验证 & \textbf{默认三档均测}：\texttt{run.sh} 依次 \texttt{aiv=1}、\texttt{aiv=2}、\texttt{aiv=8}；亦可 \texttt{--aiv 1|2|8} 单档。 \\
融合算子 AIV 约束 & 含 Cube 的融合内核（如 \texttt{sepolyvec8\_ntt\_k4}）在 910B 上 Vector 数须为 \textbf{2 的倍数}（典型 MIX 1 AIC : 2 AIV）。本用例纯 AIV，增测 \texttt{aiv=2} 以验证\textbf{最小偶数多核}分核与流水，便于与融合 launch 对齐。 \\
验证命令 & Python 先生成 \texttt{output/golden.bin}；每档 launch 后 \texttt{md5sum} + \texttt{cmp} + \texttt{verify\_result.py}。 \\
\bottomrule
\end{longtable}

\textbf{选型理由}：
\begin{itemize}
  \item \textbf{aiv=8}：$K=8$ 按 poly embarrassingly parallel，满核数据并行。
  \item \textbf{aiv=2}：每核 4 poly；覆盖融合路径要求的\textbf{偶数 AIV} 最小规模，验证分核偏移与多 poly 串行。
  \item \textbf{aiv=1}：单核 golden 基线，便于 Stage 隔离联调。
  \item 核入口：\texttt{blockNum==1} 串行 8 poly；否则 \texttt{polysPerBlock}$=8/\texttt{blockNum}$，\texttt{startPoly}$=\texttt{blockIdx}\cdot\texttt{polysPerBlock}$，核内逐 poly 新建 op。
\end{itemize}

\subsection{并行与分块}
\begin{itemize}
  \item 核规模见 §\ref{sec:aicore}；每 poly 仍为 256 系数、\texttt{kTileLength=32} 流水。
  \item \texttt{TPipe + TQue} 双缓冲：\texttt{kTileNum=4}，\texttt{kBufferNum=2} $\Rightarrow$ \texttt{kTileLength=32}。
  \item 流水：\texttt{CopyIn} $\rightarrow$ \texttt{Compute} $\rightarrow$ \texttt{CopyOut}，与 \texttt{AddCustomSample} 同构。
\end{itemize}

\subsection{计算步骤与 API 映射}
\begin{enumerate}
  \item \textbf{GM$\rightarrow$UB}：\texttt{DataCopy} 读 \texttt{se} 一 tile（\texttt{int32}）。
  \item \textbf{求 $hi$}：\texttt{ShiftRight(hi, v, 6)}。
  \item \textbf{求 $lo$}：\texttt{Muls(tmp, hi, 64)}；\texttt{Sub(lo, v, tmp)}。
  \item \textbf{窄化 int8}：\texttt{CastI32ToI8} 链（§\ref{sec:cast}）。
  \item \textbf{UB$\rightarrow$GM}：\texttt{DataCopy} 写 HI/LO 行。
\end{enumerate}

\section{全量 API 清单（查阅 CANN 9.0.0）}
\label{sec:api-table}

说明：\textbf{分类}按 PDF 第 2 章；\textbf{在线页}为社区 \texttt{atlasascendc\_api\_07\_*}；\textbf{PDF 约页}据目录（2.3.1 自 p.136，2.3.3 自 p.357，具体 API 页码因文档排版略有偏差，以在线页为准）。

\begin{longtable}{@{}L{2.6cm}L{1.5cm}L{2.4cm}L{0.8cm}L{3.6cm}L{3.5cm}@{}}
\toprule
API & 分类 & 在线文档 ID & A2 & 输入/输出 dtype（本用例） & 备注 \\
\midrule
\endhead
\apiname{GlobalTensor} & 2.2 数据结构 & 07\_0007 & $\checkmark$ & GM 张量描述 & \texttt{SetGlobalBuffer} \\
\apiname{LocalTensor} & 2.2 数据结构 & 07\_0006 & $\checkmark$ & UB 张量描述 & \texttt{AllocTensor}/\texttt{Get} \\
\apiname{GetBlockIdx} & 2.3 系统变量 & 07\_0185 & $\checkmark$ & 返回核索引 & \texttt{int64\_t} \\
\apiname{KERNEL\_TASK\_TYPE\_DEFAULT} & Utils / 核类型 & 见编程指南 & $\checkmark$ & \texttt{KERNEL\_TYPE\_AIV\_ONLY} & 纯向量核 \\
\apiname{TPipe} & 2.3 资源管理 & 07\_0108 & $\checkmark$ & 管道对象 & 每 kernel 唯一实例 \\
\apiname{InitBuffer} & 2.3 资源管理 & 07\_0110 & $\checkmark$ & 为 \texttt{TQue}/\texttt{TBuf} 分配 UB & \texttt{TPipe} 成员 \\
\apiname{TQue} & 2.3 Pipe/Que & 07\_0137 & $\checkmark$ & 队列；\texttt{VECIN}/\texttt{VECOUT} & \texttt{EnQue}/\texttt{DeQue} \\
\apiname{TBuf} & 2.3 Pipe/Que & 07\_0137 同章 & $\checkmark$ & \texttt{VECCALC} 临时区 & \texttt{Get}/\texttt{GetWithOffset} \\
\apiname{DataCopy} & 2.3.1 搬运 & 07\_0101 & $\checkmark$ & \texttt{int32}/\texttt{int8} GM$\leftrightarrow$UB & 连续块；本用例对齐 \\
\apiname{ShiftRight} & 2.3.3 矢量 & 07\_0059 & $\checkmark$ & src/dst：\texttt{int32\_t}；标量 $[0,32]$ & 算术右移 \\
\apiname{Muls} & 2.3.3 矢量 & 07\_0055 & $\checkmark$ & src/dst：\texttt{int32\_t}；标量 \texttt{int32} & $hi\times 64$ \\
\apiname{Sub} & 2.3.3 矢量 & 07\_0036 & $\checkmark$ & src/dst：\texttt{int32\_t} & $v-hi\cdot64$ \\
\apiname{Cast} & 2.3.3 类型转换 & 07\_0073 & $\checkmark$ & 见 §\ref{sec:cast} & 分步窄化 \\
\bottomrule
\end{longtable}

\subsection{曾评估但未采用的 API}
\label{sec:rejected}
\begin{longtable}{@{}L{2.6cm}L{2.2cm}L{0.8cm}L{3.8cm}L{4.8cm}@{}}
\toprule
API & 在线 ID & A2 & 原计划 & 结论 \\
\midrule
\apiname{And} & 07\_0062 附近 & $\checkmark$ & $lo=v\land 63$ & \texttt{int16\_t}/\texttt{uint16\_t}；\textbf{非 int32}。且 CPU 路径曾未得到正确 $lo$ \\
\apiname{Cast(int32$\rightarrow$int8)} & 07\_0073 & $\times$ & 直接窄化 & A2 表 6 无此组合；dav\_c220 \texttt{CheckCastDatatype} 拒绝 \\
\apiname{Cast(int16$\rightarrow$int8)} & 07\_0073 & $\times$ & 二级窄化 & 同上；CPU 仿真 \texttt{abort} \\
\bottomrule
\end{longtable}

\section{Cast 窄化链与备选}
\label{sec:cast}

\subsection{目标}
将 $hi, lo \in [0,63]$ 的 \texttt{int32} 写入 \texttt{int8} 输出（值域安全，无饱和风险）。

\subsection{A2 文档表 6 允许的组合（摘录）}
\begin{itemize}
  \item \texttt{int32\_t} $\rightarrow$ \texttt{int16\_t}：\texttt{CAST\_NONE}
  \item \texttt{int16\_t} $\rightarrow$ \texttt{half}：\texttt{CAST\_NONE}（或与 \texttt{SetDeqScale} 配合）
  \item \texttt{half} $\rightarrow$ \texttt{int8\_t}：多种 \texttt{RoundMode}
\end{itemize}

\subsection{实现链（已验证）}
\begin{verbatim}
Cast(tmpI16, srcI32, CAST_NONE, count);
Cast(tmpHalf, tmpI16, CAST_NONE, count);
Cast(dstI8, tmpHalf, CAST_NONE, count);
\end{verbatim}

\subsection{备选方案（若 half 链不可用）}
\begin{enumerate}
  \item \textbf{备选 A}：\texttt{int32}$\rightarrow$\texttt{float}$\rightarrow$\texttt{int8}（需核对表 6 是否含 \texttt{int32}$\rightarrow$\texttt{float} 与 \texttt{float}$\rightarrow$\texttt{int8}；有舍入语义）。
  \item \textbf{备选 B}：标量循环写 \texttt{int8}（仅 32 元/tile，可接受但失去向量窄化）。
  \item \textbf{备选 C}：查 \texttt{MulCast} 等复合 API（本用例值域小，非必须）。
\end{enumerate}

\section{逐步代入与覆盖率评估}

\begin{longtable}{@{}L{1.0cm}L{3.8cm}L{4.2cm}L{2.8cm}@{}}
\toprule
步骤 & 数学 & API & 覆盖 \\
\midrule
S0 & 核映射 & \texttt{GetBlockIdx}, \texttt{KERNEL\_TASK\_TYPE\_DEFAULT} & 100\% \\
S1 & 读 $v$ & \texttt{DataCopy} + \texttt{TQue} & 100\% \\
S2 & $hi=v\gg 6$ & \texttt{ShiftRight} & 100\% \\
S3 & $lo=v-64\cdot hi$ & \texttt{Muls}, \texttt{Sub} & 100\% \\
S4 & 写 int8 & \texttt{Cast}$\times 3$ & 100\%（链式） \\
S5 & 写 GM & \texttt{DataCopy} & 100\% \\
\midrule
\multicolumn{3}{l}{\textbf{合计（主路径）}} & \textbf{100\% 向量} \\
\bottomrule
\end{longtable}

\textbf{结论}：在采用 $lo=\texttt{Sub}(v,\texttt{Muls}(hi,64))$ 与三段 \texttt{Cast} 前提下，Stage1 主路径 \textbf{无需标量回退}。标量仅存在于核 launch 壳（\texttt{main.cpp}）与 Python Golden。

\section{源码与验证}

\begin{itemize}
  \item 内核：\texttt{f203\_stage1\_encode\_custom.cpp}
  \item Launch：\texttt{launch\_profile.h}（\texttt{aiv=1}/\texttt{aiv=8}）
  \item 验证：\texttt{bash run.sh -r cpu -v Ascend910B4}（默认两档；golden \texttt{md5sum} 对拍）
  \item 索引：\texttt{library/documents/CANN-AscendC算子开发接口参考-查阅索引.md}（2026-05-19 条目）
\end{itemize}

\end{document}
